\section{Parâmetros arquiteturais}

\subsection{Largura física}

Cada registrador vetorial físico tem largura:

|$$
VLEN = 128 \text{ bits} = 16 \text{ bytes}.
$$|

O banco contém 32 registradores físicos, indexados de 0 a 31.

\subsection{Largura do elemento}

A largura do elemento selecionado, denominada SEW, pode assumir os
valores suportados pelo modelo:

|$$
SEW \in \{8,16,32\} \text{ bits}.
$$|

O maior elemento suportado é:

|$$
ELEN = 32 \text{ bits}.
$$|

Nesta etapa, os valores são tratados como inteiros sem sinal para
representação e como inteiros de largura fixa para as operações
aritméticas. A aritmética modular é definida pela largura do elemento.

\subsection{LMUL}

A configuração LMUL suportada nesta versão é:

|$$
LMUL \in \{1,2,4,8\}.
$$|

A implementação não inclui multiplicadores fracionários de registradores,
como LMUL igual a |$1/2$|, |$1/4$| ou |$1/8$|.

\subsection{VLMAX}

Para os valores inteiros de LMUL implementados, o número máximo de
elementos ativos em um grupo vetorial é:

|$$
VLMAX = \frac{VLEN}{SEW} \times LMUL.
$$|

Exemplos:

\begin{center}
\begin{tabular}{ccc}
\toprule
SEW & LMUL & VLMAX \\
\midrule
8  & 1 & 16 \\
16 & 1 & 8 \\
32 & 1 & 4 \\
32 & 2 & 8 \\
8  & 4 & 64 \\
\bottomrule
\end{tabular}
\end{center}

O comprimento efetivo do vetor, \texttt{vl}, é limitado a VLMAX:

|$$
vl = \min(vl_{\text{requested}}, VLMAX).
$$|

O tratamento de valores negativos para \texttt{vl} ainda deve ser
formalizado e validado explicitamente.